%% kapa: exam=1 round=29 year=2018 period=2 form=A subject=appraisal_law subject_ko=감정평가관계법규 seq=6 number=6
%% source: Q-Net 감정평가사 기출문제 (cst00309) artlSeq=5202808 file=extracted/1008/5202808/2018년 제29회 감정평가사 1차시험 기출문제/2018년 제29회 감정평가사 1차시험 기출문제/2교시 A형.pdf
%% conversion: kordoc 4.24.51 -> tools/gen1.py (auto); review=proofread-sonnet; flags=none
\begin{question}{6}
국토의 계획 및 이용에 관한 법령상 공동구에 관한 설명으로 옳은 것은?
\begin{choices}
  \item 「도시개발법」에 따른 도시개발구역이 100만 $\mathrm{m}^2$를 초과하는 경우에는 해당 구역에서 개발사업을 시행하는 자는 공동구를 설치하여야 한다.
  \item 공동구가 설치된 경우에 전선로, 통신선로 및 수도관은 공동구에 수용하지 아니할 수 있다.
  \item 공동구 설치에 필요한 비용은「국토의 계획 및 이용에 관한 법률」이나 다른 법률에 특별한 규정이 있는 경우를 제외하고는 사업시행자가 단독으로 부담한다.
  \item 공동구의 관리에 소요되는 비용은 그 공동구를 점용하는 자가 함께 부담하되, 부담비율은 점용면적을 고려하여 공동구관리자가 정한다.
  \item 공동구관리자는 3년마다 해당 공동구의 안전 및 유지관리계획을 수립ㆍ시행하여야 한다.
\end{choices}
\end{question}
\answer{4}
